\subsection{环的示例}

I. 零环。设 A 为一个环。在 A 中 0 = 1 的必要且充分条件是 A 由单个元素组成。该条件显然充分。另一方面，若 0 = 1，则对所有 \(x \in A\) ，x = x.1 = x.0 = 0。这样的环称为零环。

\begin{enumerate}
\def\labelenumi{\Roman{enumi}.}
\setcounter{enumi}{1}
\tightlist
\item
  有理整数环。利用 § 2 第 5 号中定义的加法，以及 § 2 第 6 号中定义的乘法，Z 是一个交换环。记号 0、1、-x 与先前引入的记号一致。
\end{enumerate}

*III. 实值函数环。设 I 为实数集 R 中的一个区间，并设 A 为定义在 I 上且取实值的连续函数所构成的集合。两个函数 f 和 g 的和 \(f + g\) 与乘积 f·g 定义为

\[
(f + g) (t) = f (t) + g (t), \quad (f g) (t) = f (t) g (t) \quad (t \in \mathbf {I}).
\]

由此得到一个交换环，其单位元为常数 1。*

*IV. 卷积伪环。设 E 为 R 上在某个有界区间之外为零的实值连续函数的集合。两个函数的和按 III 中那样定义，但乘积现在定义为

\[
(f g) (t) = \int_ {- \infty} ^ {\infty} f (s) g (t - s) d s
\]

（“卷积乘积”）。由此得到一个交换伪环，它不是环（参见《积分》VIII，§ 4）。*

V. 环 A 的反环。设 A 为一个环。带有与 A 相同的加法以及乘法 \((x,y)\mapsto yx\) 的集合 A 通常记作 \(A^{0}\) 。它是一个环（称为 A 的反环），与 A 具有相同的零元和单位元，并且当且仅当 A 为交换环时它与 A 相同。

\begin{enumerate}
\def\labelenumi{\Roman{enumi}.}
\setcounter{enumi}{5}
\tightlist
\item
  交换群的自同态环。设 G 是一个写成加法形式的交换群。设 E 表示 G 的自同态组成的集合。给定 E 中的 f 和 g，映射 \(f + g\) 和 fg （从 G 到 G 中）定义为
\end{enumerate}

\[
(f + g) (x) = f (x) + g (x), \quad (f g) (x) = f (g (x)) \quad (x \in \mathbf {G}).
\]

由 § 1 第 5 号的命题 5， \(f + g\) 是 G 的一个自同态，显然 \(fg = f \circ g\) 也是 G 的自同态。根据 § 4 第 8 号，E 在加法下构成一个（交换）群。乘法显然满足结合律且具有单位元 \(\operatorname{Id}_{\mathbf{G}}\) 。此外，对 E 中的 \(f, g\) 和 \(h\) ，设 \(\phi = f \cdot (g + h)\) ；对所有 \(x \in \mathbf{G}\) ，

\[
\phi (x) = f ((g + h) (x)) = f (g (x) + h (x)) = f (g (x)) + f (h (x))
\]

因为 \(f\) 是 \(\mathbf{G}\) 的自同态；因此 \(\phi = fg + fh\) ，并且显然

\[
(g + h) f = g f + h f.
\]

因此 E 是一个（一般不交换的）环，称为 G 的自同态环。

\begin{enumerate}
\def\labelenumi{\Roman{enumi}.}
\setcounter{enumi}{6}
\tightlist
\item
  零平方伪环。称伪环 A 为零平方的，如果对所有 \(x, y \in A\) 都有 xy = 0。设 G 是一个交换群。若给集合 G 赋予群 G 的加法以及乘法 \((x, y) \mapsto 0\) ，则得到一个零平方的伪环。它仅在 \(G = \{0\}\) 时是环，此时它是零环。
\end{enumerate}

